\section{Protobuf Compatibility}\label{sec:proto-compat}

This section will detail the Lean formalization of Protobuf that's being
developed for Pollux.

\subsection{Descriptor Design}

The intermediate encoding format used a custom extensionality theorem over a
descriptor structured as a list to show that descriptors were equal. This was
required since Lean's finite maps (unlike Rocq's \texttt{gmap}) can't be placed
inside of inductive types. That's a clunky choice which was not what I wanted
when I told Aristotle ``I want extensionality on my descriptors'', but I kinda
just lived with it. Well, the full protobuf descriptors we're taking a more
principled approach that can fall back on the vast support of Mathlib.


\subsubsection{Protobuf Descriptors in Lean}

Instead, the full protobuf descriptors take a different approach, using a
dependently typed list internally but exposing a view of the descriptor using a
proper finite map.

\begin{minted}{lean}
inductive ScalarType where
  | double | float | int32 | int64 | uint32 | uint64
  | sint32 | sint64 | fixed32 | fixed64 | sfixed32 | sfixed64
  | bool | string | bytes

inductive Cardinality where
  | singular | optional | repeated | oneof (group : Nat)

mutual
inductive Desc      where | mk (es : List ((_ : Int) × Field))
inductive Field     where | mk (card : Cardinality) (ty : FieldType)
inductive FieldType where | scalar (s : ScalarType) | msg (d : Desc)
end
\end{minted}

Outside of Lean, we'll use this syntax for the Protobuf descriptors.

\begin{figure}[h]
  \centering
  \begin{syntax}
    \categoryFromSet[Naturals]{n}{\mathbb{N}}

    \category[Scalar Type]{\tau_s}
    \alternative{\double}
    \alternative{\float}
    \alternative{\ints}
    \alternative{\intl}
    \groupleft{
        \alternative{\uints}
        \alternative{\uintl}
        \alternative{\sints}
        \alternative{\sintl}
    }
    \groupleft{
      \alternative{\fixeds}
      \alternative{\fixedl}
      \alternative{\sfixeds}
      \alternative{\sfixedl}
    }
    \groupleft{
      \alternative{\bool}
      \alternative{\str}
      \alternative{\byt}
    }

    \category[Cardinality]{c}
    \alternative{\sig}
    \alternative{\opt}
    \alternative{\rep}
    \alternative{\oneof{n}}

    \category[Field]{f}
    \alternative{\langle c, \tau_s \rangle}
    \alternative{\langle c, d \rangle}

    \category[Field List]{l}
    \alternative{\texttt{[]}}
    \alternative{(n, f) :: l}

    \category[Reserved Set]{r}
    \alternative{\varnothing_r}
    \alternative{n, r}

    \category[Descriptor]{d}
    \alternative{\llangle r, l \rrangle}
  \end{syntax}
  \caption{Pollux syntax for Protobuf descriptors.}\label{fig:proto-desc-syn}
\end{figure}

The use of the sigma type is important, since it let's us tie into the existing
Mathlib infrastructure for building a finite map from from the list
internals. The basic idea is that the descriptor itself is \emph{sealed}, and we
can explode one layer at a time to get a view of the descriptor through the lens
of a finite map which should enable reasoning about the descriptor using the
theory from Mathlib. This approach also is designed to respect these
requirements:

\begin{itemize}
\item \textbf{Kernel Acceptance.} Clearly we need a representation that Lean can
  handle. Under this choice, we're not storing the map directly but are still
  able to indirectly access other theorems about it.
\item \textbf{Extensionality.} The existing formulation certain variants of the
  compatibility theorems state that we recover not an equivalent value, but an
  \emph{equal} value, which \texttt{Finmap} provides.
\item \textbf{Termination Measures.} Since we have nested messages, we need to
  ensure that we can safely traverse the descriptor and recurse on the
  descriptor. Pervious formats, like the intermediate one, defined a custom
  termination metric by noting that the depth nested descriptor must be smaller
  than the containing descriptor. This works well for tree-like descriptors, but
  Protobuf does support recursive and mutually-recursive messages definitions,
  which cannot be represented by this descriptor type.
\item \textbf{Forward Compatibility.} Protobuf can represent it's own
  descriptors, but not without recursive messages. Eventually we'll have to
  support getting a descriptor into the Lean formalization, and one option is to
  model a descriptor as a value of \texttt{descriptor.proto} and then build the
  descriptor from that information. This is common in real-world protobuf
  implementations, and the sealed design should still work for this approach.
\end{itemize}

There is a well-formedness requirement on descriptors, one coming from the
internal structure of the descriptor in Lean (Definition~\ref{def:dwf}) and the
other from the Protobuf compiler (Definition~\ref{def:dlegal}).

\vspace{1em}
\begin{definition}[Descriptor Well-Formedness]\label{def:dwf}
  A descriptor $d$ is well-formed, denoted $d \checkmark_w$, if all fields in
  the descriptor have a unique field number, the fields are stored in sorted
  order, and all nested descriptors are also well-formed.
\end{definition}

\vspace{1em}
\begin{definition}[Legal Descriptor]\label{def:dlegal}
  A descriptor $d$ is legal, denoted $d \checkmark_l$, if
  \begin{itemize}
  \item All field numbers are between 1 and $2^{29}-1$, excluding the range
    19,000 to 19,999, which are reserved by the Protobuf compiler.
  \item Message fields are always explicitly marked as optional. This conforms
    with Protobuf, which will generate a method to check message presence, and
    will allow us to initialize recursive messages.
  \item All nested descriptors within $d$ are also legal.
  \end{itemize}
\end{definition}

In preparation for needing to support recursive messages, Pollux uses
well-founded recursion on the size of the descriptor through the \texttt{get?}
interface. We define a function to compute the size of the descriptor below.

\vspace{1em}
\begin{definition}[Descriptor Size]\label{def:dsize}
  The size of a descriptor $d$ is defined as $\mathrm{size}_d$ below.
  \begin{mathpar}

    \mathrm{size}_d\left(\llangle r, l \rrangle\right) = 1 + \mathrm{size}_d(l)

    \mathrm{size}_d(\langle \_, \tau_s \rangle) = 1
    
    \mathrm{size}_d(\langle \_, d \rangle) = 1 + \mathrm{size}_d(d) 

    \mathrm{size}_d(l) =
    \begin{cases}
      0 & \text{if } l = \texttt{[]} \\
      \mathrm{size}_d(f) + \mathrm{size}_d(l') & \text{if } l = (\_, f) :: l'
    \end{cases}
  \end{mathpar}
\end{definition}

Finally, the \texttt{oneof} cardinality needs some explanation. When looking at
the syntactic structure of a Protobuf message in a \texttt{.proto} file, the
logical way to model a \texttt{oneof} is likely to have a set of fields or some
other similar grouped structure. However, this isn't how Protobuf actually
models a \texttt{oneof} internally, and there are good reasons for Pollux to
follow their lead. 

\begin{itemize}
\item Members of a \texttt{oneof} share a field number namespace with the rest
  of the message. You can't define a field with identifier $n$ in two
  \texttt{oneof}'s in the same message, and using the \texttt{oneof} cardinality
  like we have it lets the well-formedness condition handle this automatically.
\item A \texttt{oneof} field doesn't have a unique field number, so there is no
  key available to insert a dedicated \texttt{oneof} representation into the
  finite map the descriptor is modeled as.
\item A field in a \texttt{oneof} can't be repeated, and \texttt{oneof}'s can't
  be empty. Both of those requirements are structurally impossible.
\item The Protobuf underlying \texttt{oneof} index isn't evolution stable, since
  introducing a new \texttt{oneof} will shift the assigned index of everything
  later. This complicates the \texttt{oneof} compatibility rules, but at least
  by adopting the same representation, everything is transparent.
\item If a program called \texttt{get?} with a field in a \texttt{oneof}, the
  \texttt{get?} operation would need to explore the inside of the \texttt{oneof}
  (via a mechanism perhaps similar to \texttt{explode}) anyways.
\end{itemize}

Fortunately, checking equality on the natural number included with each
\texttt{oneof} cardinality provides the required information. Even with a nested
or grouped representation, the parser would still have to manually unset the
previously set field, so this representation doesn't add any unnecessary work.

\subsubsection{Sealed \& Unsealed}

While the descriptor internals use a sorted listed, the public interface should
respect a \emph{seal} by convention (lean doesn't have great mechanisms to mark
things private or opaque). The public interface is this: 

\begin{minted}{lean}
def explode (d : Desc) : Finmap (fun _ : Int => Field) :=
    d.entries.toFinmap
def get?    (d : Desc) (k : Int) : Option Field := d.explode.lookup k
\end{minted}

Where \texttt{explode} unwraps the outer-most descriptor into a Mathlib
\texttt{Finmap}, enabling proofs to use the established \texttt{Finmap}
theory. Because \texttt{Finmap} quotients over the ordering of values in the
map, the sorted-list invariant is invisible most consumers of the interface, as
long as they use the provided mechanisms when inserting into a descriptor. Some
operations do still require the well-formed predicate, such as extensional
equality, deletions and a helpful lemma which shows that if two descriptors have
the same exploded result, they must be the same descriptor.

Descriptors should be sealed. Looking at how the code for the intermediate
format is structured, descriptors are never consumed more than one layer at a
time. The serializer iterates values, not the descriptor. While looking at a
message, we check against fields as we encounter them in the value.

However, values do need to iterated over and are the basis for round-trip
proofs. Sealing the values within the parser and serializer implementation would
create more work than we desire, although the theorems will still be able to use
a variant of the explode operator if it makes more sense.

\subsubsection{Considerations about Ordering \& Determinism}

The Protobuf spec permits serializing and parsing fields in any order, which
seems at odd with the choice to represent the fields in a sorted list. It should
be possible to explode a layer and use the undefined iteration order of the
\texttt{Finmap} to be as generalize as possible with regards to that Protobuf
spec. However,

\begin{itemize}
\item Having the serializer return data in any order would mean that it isn't a
  pure function, which is difficult to represent in Lean. We could quotient over
  the ordering of the fields, but then we'd have to prove that every other
  function consuming the serialized output also respects the quotient type,
  which is possible but a lot of work.
\item The exact output from the parser is descriptor dependent as well as value
  dependent (recall that we can't know what value is encoded without needing to
  to a descriptor, see Section~\ref{sec:philosophy}). Having both align as an
  ordered list reduces map equivalences to syntactic list equality.
\item The Protobuf spec is also adding a requirement to the parser, that is has
  to be able to handle encodings which were produced with the fields in an
  arbitrary order, with last field wins semantics. The existing \texttt{ParseOk}
  definition used for the intermediate formalization only requires the parser to
  be correct on the output of our own serializer, not any possible encoding. It
  will need to be extended with a relational specification which allows the
  maximum possible amount of non-determinism. The soundness requirement for this
  specification would require that our serializer produces values that satisfy
  the relational spec while completeness would require that any encoding in the
  relational specification can be parsed correctly.
\end{itemize}

\subsubsection{New Protobuf Features}

There are three new, notable features in Protobuf that were missing from the
intermediate formate entirely.

\begin{itemize}
\item \textbf{Cardinality.} The intermediate format only had scalar values, but
  full Protobuf has repeated values. (Optional values and implicit values are
  already extremely similar, so the addition of repeated values is the most
  important). 
\item \textbf{Maps.} Maps are easy to support since we can just desugar them, as
  described in Section~\ref{sec:proto}, with a few additions for the
  restrictions on map key type and the like.
\item \textbf{Oneof.} On the wire, a \texttt{oneof} is invisible. There is no
  information included in the wire level format to describe that a field is in a
  \texttt{oneof}, it's another one of those things that we have to have a
  descriptor to know about. Individual fields still have their own field number,
  which is encoded, that comes from the name space of the whole message. The
  \texttt{oneof} really adds a parser semantic requirement about clearing the
  old field if we encounter a new encoded field in the same \texttt{oneof}, and
  overall constraint on the final message that at most one of the fields will be
  set. The \texttt{oneof} is then represented as a cardinality, with some extra
  information to determine which \texttt{oneof} in the message this field
  belongs to. This exactly mirrors how the \texttt{descriptor.proto} operates,
  and we can use natural number equality to scan the value and clear old set
  fields when needed. Since each field has exactly one cardinality, this already
  enforces that we can't have repeated fields inside a oneof and that a oneof
  can't be empty.
\end{itemize}

\subsubsection{Omitted Protobuf Features}

At the moment, these features are \emph{not} modeled in Pollux. Adding support
for them should be relatively straight-forward. If I've already developed
relevant rules for them, they might be included in a different color.

\textbf{Some of these features will likely be added before the work is complete.}

\begin{itemize}
\item \textbf{Reserved field numbers.} Protobuf doesn't recommend reusing field
  numbers for new fields (i.e.\ removing a field, then later adding a field with
  that same number), and that also poses a challenge to Pollux. If a field were
  added which uses and old field number, Pollux would assume that the intent is
  for the new field to be compatible with the old one. Since we expect the
  compatibility relation to be transitive, any evolving derivation of that
  message would attempt to force a connection between two fields that are
  semantically unrelated.

  Fortunately, Protobuf has a mechanism for this, reserved field numbers. You
  can set specific field numbers as reserved, then the protobuf compiler will
  ensure that you don't add a field with that number to the message in the
  future. Pollux will also enforce this, even though the Lean formalization
  doesn't track it \emph{yet}.

\item \textbf{Field names.} Every Protobuf field has a name, which is used in
  the generated code from \texttt{protoc}. These names aren't recorded at all in
  the wire format, but do impact the compatibility of a change with the rest of
  your code. At the moment, that's considered beyond the scope of Pollux.

\item \textbf{Reserved field names.} Like reserved field numbers, you can also
  reserve field names. Pollux's stance is that this is outside the domain we're
  targeting, but could be added with mechanisms similar to the reserved field
  number.

\item \textbf{Recursive Messages.} Protobuf messages (and thus descriptors) can
  be recursive. The current lean formalization enforces a tree structure on the
  descriptors. This is nice for recursion over the descriptors, but does mean we
  can't model all descriptors. The Protobuf data set does contain recursive and
  mutually recursive descriptors. It seems that these are more common in tools
  which analyze Protobuf descriptors, likely because \texttt{descriptor.proto},
  which models Protobuf descriptors as a Protobuf struct, is recursive. If we
  want to transmit protobuf descriptors into the Pollux lean formalization using
  \texttt{descriptor.proto} (a reasonable choice), we'll have to expand the
  descriptor representation to be a flat-map of named messages. See
  Table~\ref{tab:recur-summary} and Table~\ref{tab:recur-repos} for information
  about recursive messages in the Protobuf data set.

\item \textbf{Enums.} Protobuf supports enums, which shouldn't be hard to
  support in the future, but aren't present currently.

\item \textbf{Groups.} This is a Protobuf 2 feature, which still has some
  reserved encoding wire type implications. It's simply not used in Protobuf 3.

\item \textbf{Explicit Maps.} Maps can be desugared to a \texttt{repeated}
  key-value pair (See Section~\ref{sec:proto}), so no explicit representation is
  needed. Ok, this is technically \emph{not} true, since maps have additional
  requirements about the type of the key, but if we were to desugar a map to
  it's internal representation, the user would need to also desugar it to create
  something which might be problematic.

\item \textbf{Packed Scalars.} Not really a concern of the descriptor, although
  there are serious consequences of this.
\end{itemize}

\begin{table}[H]
\centering
\begin{tabular}{lcc}
    \toprule
    Metric & Count & Percentage \\
    \midrule
    Messages & 57,461 & \\
    Recursive & 831 & 1.45\% \\
    \bullet{} self-recursive & 425 & 0.74\% \\
    \bullet{} mutually-recursive & 447 & 0.78\% \\
    Contains recursion & 3,893 & 6.78\% \\
    \bottomrule
\end{tabular}
\caption[Summary of recursive protobuf messages.]{Summary statistics of
    recursive descriptor usage in the Protobuf data set. Here ``recursive''
    means directly self-recursive or in a mutually recursive cycle while
    ``contains recursion'' is any message which includes a recursive message in
    it's transitive closure. Any message which contains recursive cannot be
    represented by the current descriptor model. There are 533
    mutually-recursive cycles with the longest containing 22 messages. Note
    that there are 41 messages which are both self-recursive and
    mutually-recursive.}\label{tab:recur-summary}
\end{table}

\begin{table}[H]
\centering
\begin{tabular}{lccccc}
    \toprule
    Repository & Messages & Recursive & \% & Contains & \% \\
    \midrule
    \texttt{protocolbuffers/protobuf-go} & 605 & 110 & 18.18\% & 151 & 24.96\% \\
    \texttt{cloudwego/kitex} & 34 & 4 & 11.76\% & 7 & 20.59\% \\
    \texttt{go-kratos/kratos} & 94 & 8 & 8.51\% & 17 & 18.09\% \\
    \texttt{bufbuild/buf} & 2,922 & 210 & 7.19\% & 371 & 12.70\% \\
    \texttt{google/protobuf.dart} & 503 & 30 & 5.96\% & 59 & 11.73\% \\
    \texttt{protocolbuffers/protobuf} & 2,372 & 172 & 7.25\% & 242 & 10.20\% \\
    \texttt{googleapis/googleapis} & 49,370 & 289 & 0.59\% & 3,022 & 6.12\% \\
    \texttt{streamdal/plumber} & 286 & 2 & 0.70\% & 13 & 4.55\% \\
    \texttt{google/rejoiner} & 50 & 0 & 0.00\% & 2 & 4.00\% \\
    \texttt{googleapis/google-cloud-go} & 134 & 3 & 2.24\% & 3 & 2.24\% \\
    \texttt{uber/prototool} & 469 & 3 & 0.64\% & 6 & 1.28\% \\
    \texttt{asynkron/protactor-go} & 149 & 0 & 0.00\% & 0 & 0.00\% \\
    \texttt{colinmarc/hdfs} & 443 & 0 & 0.00\% & 0 & 0.00\% \\
    \texttt{twitchtv/twirp} & 30 & 0 & 0.00\% & 0 & 0.00\% \\
    \bottomrule
\end{tabular}
\caption[Per-repo recursive message breakdown]{A pre-repo breakdown of
    recursive message and messages which contain recursion. Notice that tools
    which work with protobuf descriptors have a higher percentage of recursive
    messages than project which just use Protobuf. The aggregate numbers in
    Table~\ref{tab:recur-summary} are dominated by
    \texttt{googleapis/googleapis} since it contains 86\% of the data set. I
    don't know what's happening with \texttt{google/rejoiner}, but it's
    possible some messages didn't compile correctly.}\label{tab:recur-repos}
\end{table}

\subsection{Value Design}

Mirroring the descriptors are the values, which carry actual data rather than
describe the shape of a message.

\subsubsection{Values in Lean}

In the Lean formalization, the structure of a value matches the structure of
descriptors across all three levels.

\begin{itemize}
\item A descriptor can be matched to a value.
\item A field can be matched to a val. {\color{red} I'd like to change the
    name.}
\item A field type can be matched to a payload.
\end{itemize}

\begin{minted}{lean}
mutual
inductive Value   where | mk (es : List ((_ : Int) × Val))
inductive Val     where | implicit (p : Payload)
                        | optional (p : Option Payload)
                        | repeated (ps : List Payload)
inductive Payload where | int (z : Int) | bool (b : Bool) | string (s : String)
                        | bytes (bs : List UInt8) | float (bits : UInt32)
                        | double (bits : UInt64) | msg (v : Value)
end
\end{minted}

Notice that just like with the descriptor, the presence wrapper sites above the
payload, and also reduces the number of constructor cases (marginally) and
creates a regular structure for \texttt{optional} and \texttt{repeated} fields
of an \texttt{option} / \texttt{list} of payloads rather than an \texttt{option}
/ \texttt{list} with a specific instance type. Actual data is represented in
Lean's native types. The \texttt{msg} payload constructor is the only major
structural difference from the descriptor, since it sits in the \texttt{Payload}
type rather than in a separate mutually inductive type. (I've eliminated some of
the difference in the formal grammar, but the distinction on the lean side is
worth noting.) However, the inductive structure \emph{does} match --- it's
under the presence modifier, but still above the leaf of the descriptor or value
--- and the descriptor's scalar type gains derivable equality by being outside the
mutually recursive block. The way that the value \texttt{Payload}'s are used,
this isn't required, they're compared against a field from the descriptor in a
way to avoid this problem. Also, doubles and floats are represented as unsigned
integers to get bit-level equality. Lean's formalization of IEEE 754 floating
point numbers is opaque and the defined equality isn't reflexive (NaN $\ne$ NaN). 

{\color{red} This might change. It's possible to pull the scalar payloads out if
needed.}

\begin{figure}[h]
  \centering
  \begin{syntax}
    \categoryFromSet[Integers]{z}{\mathbb{Z}}

    \categoryFromSet[Booleans]{b}{\mathbb{B}}

    \abstractCategory[Byte Strings]{bs}

    \abstractCategory[Strings]{s}

    \categoryFromSet[32-bit Unsigned Integers]{\mathrm{u^{32}}}{[0, 2^{32})}

    \categoryFromSet[64-bit Unsigned Integers]{\mathrm{u^{64}}}{[0, 2^{64})}

    \category[Payload]{p}
    \alternative{z}
    \alternative{b}
    \alternative{s}
    \alternative{bs}
    \alternative{\mathrm{u^{32}}}
    \alternative{\mathrm{u^{64}}}
    \alternative{v}

    \category[Val]{w}
    \alternative{\impt{p}}
    \alternative{\optt{\text{none}}}
    \alternative{\optt{p}}
    \alternative{\rept{[p_1; \dots; p_m]}}

    \category[Value List]{l}
    \alternative{\texttt{[]}}
    \alternative{(n, v) :: l}

    \category[Value]{v}
    \alternative{\llangle l \rrangle_v}
  \end{syntax}
  \caption{Pollux syntax for Protobuf values.}\label{fig:proto-val-syn}
\end{figure}

\subsubsection{Value Totality}\label{sec:value-totality}

In real-world Protobuf implementations, a value always has the exact same domain
as the descriptor which generated it. If a message has a missing field (for any
reason) and that field is marked as implicit in the descriptor, the default
value is automatically populated into the field. Pollux matches this model by
requiring that when a value is linked to a specific descriptor, the domains of
the two match exactly. This will be part of the value well-formedness
conditions. Pollux will model a missing field using either the default value for
that type, \texttt{none} for optional values and an empty list for
\texttt{repeated} ones, just like Protobuf.

This choice strengthens the connection between a descriptor and values under it,
since rather then being able to mix and match the two arbitrarily --- something
that no real-world Protobuf implementation allows --- the assumptions made by
Pollux closely model the data assumptions of Protobuf.

However, on the proof aspect of Pollux both directions are important and worth
concerning.
\begin{itemize}
\item \textbf{Declared implies present.} A field being declared in the
  descriptor implies that something is present here in the matching value. This
  is the \emph{no omissions requirement}.
\item \textbf{Present implies declared.} Likewise a value being present in the
  message implies that it was declared in the descriptor. This is the \emph{no
    unknown values requirement}.
\end{itemize}

Starting with the first case, this requirement arises from some trouble with the
intermediate parsing format's compatibility relation, specifically the
\textsc{M-Declare} rule, where the writer's descriptor carries a field that
isn't present in the value, and the reader gets a missing value (a
intermediate format feature not present in Protobuf, this would become the
default for the type of the field).
\[
  \infer[M-Declare]{m_1 : d_1 \preceq m_2 : d_2 \\\\ k \not \in \dom{d_1} \\ k \not \in
    \dom{m_1} \\ f_1 \propto f_2}{ m_1 : d_1[k \mapsto f_1] \preceq m_2[k \mapsto \mathtt{M\_MISSING}] :
    d_2[k \mapsto f_2]}
\]
For that format, this rule was required since without it, the message
compatibility relation could not follow the ID compatibility relation which
characterizes what could happen to a value being serialized and parsed under the
same descriptor. The well-formedness condition on the values in the intermediate
format only constrained the value without regard for a descriptor.

Under the new totality requirement, this rule is vacuous. But why is removing it
a good thing?
\begin{enumerate}
\item This doesn't happen under the Protobuf data model, which would have
  inserted the default value when creating or parsing the value originally. This
  rule only exists since the well-formedness condition of the intermediate
  format was too weak, but becomes a case in every proof about the message
  compatibility relation that describes a case that Protobuf will never express.
\item Without totality, the intermediate format never had message equality when
  serializing and parsing under the same descriptor, and this requirement
  simplifies the proofs and theorem statements.
\end{enumerate}

In the second case, the no unknown fields requirement does two important
things:

\begin{enumerate}
\item The transformations output (Section~\ref{sec:proto-transform}) will always
  have a domain of $\dom(d_2)$ since we want to provide a value when viewed
  against $d_2$. So the keys that we lose when transforming a value $v$ are
  exactly $\dom{v} \setminus \dom{d_2}$. Having exact domain equality between
  $v$ and $d_1$ means that this is the same as $\dom{d_1} \setminus \dom{d_2}$
  and enables the descriptor level compatibility relation to characterize the
  value-level transformation. The drop rule in both the descriptor compatibility
  relation and the message one for Protobuf will be applied only at write-only
  descriptor key. Allowing unknown fields to exist in the value would break
  this, and a descriptor analysis wouldn't be able to meaningful reason about
  what happens at the value level.
\item Transforming a value from it's linked descriptor to that same descriptor
  should yield exactly the same value. Without this requirement, an unknown field
  would be dropped since Pollux doesn't track unknown fields explicitly
  (Protobuf actually does, for largely this reason, motivated by a potential
  man-in-the-middle where the sender and receiver both have a larger descriptor
  then the relay).
\end{enumerate}

At the byte level, the serializer will never emit a tag that the parser using
the same descriptor can't understand, so our prospective is that an unknown
field must come from a schema difference in the descriptors (or some
modification made to the value on the network; well outside our scope).

\subsubsection{Cardinality and Scalar Carriers}

The presence marker in the value layer has three constructors, while the
cardinality marker in the descriptor layer has four, with \texttt{oneof} being
the missing one. Be default in Protobuf, all \texttt{oneof} fields are treated
as \texttt{optional}, which Pollux also respects. The requirement that only one
is set is a cross-field constraint, which will live in value's legal predicate.

Additionally, while the descriptor scalar types enumerate all Protobuf type, the
value collapses several of these into a shared carrier. Specifically for all the
various integer types, the value uses a mathematical integer when the Protobuf
scalar types change in encoding and range. The legal predicate for values will
enforce the range requirement while the encoding choice is only relevant
during the serialization and parsing process, not the final values.

\subsubsection{Value Initialization}

Since messages have to match their descriptor, it's important that we're able to
generate initial, empty messages to establish this invariant. Notice that the
\texttt{Field.init} function takes the cardinality information in addition to the
type information. This is what lets the initialization process differentiate
between the default value of a repeated field (an empty list), optional field
(\texttt{none}) and scalar field.

\begin{minted}{lean}
def Field.init : Field → Val
  | .mk .singular (.scalar s) => .implicit (ScalarType.defaultPayload s)
  | .mk .singular (.msg _)    => .optional none   -- unreachable when Legal
  | .mk .optional _           => .optional none
  | .mk .repeated _           => .repeated []
  | .mk (.oneof _) _          => .optional none

def Value.init (d : Desc) : Value := ⟨d.entries.map
    (fun e => ⟨e.1, Field.init e.2⟩)⟩

theorem Value.get?_init (d : Desc) (k : Int) :
    (Value.init d).get? k = (d.get? k).map Field.init
\end{minted}

This is one of two places where the descriptor seal is broken and (recursive) iteration over
the descriptor is performed, the other being the injection of missing fields
during parsing. 

\subsubsection{Value Validity} 

Descriptor validity broke down into two requirements, a Lean specific invariant
on the internal structure of our representation, and the Legal predicate that
enforces some particular requirements of the Protobuf compiler. A value being
well-formed is a bit more complex since we require that it's well-formed
\emph{against} a descriptor. The value well-formedness predicate will be used as
the output well-formedness predicate of the Protobuf serializer, meaning we'll
be able to assume that a value input to the serializer is well-formed and
use that in combination progress theorem that all well-formed values will
serialize. Validity on the value is the precondition for the top-level
round-trip theorems, and are encoded into the type of the serializer from our
serializer combinator library. The actual contents of the validity predicate are
stronger than needed to show that the serialization will always succeed, but can
be weakened to imply the needed combinator well-formedness conditions throughout the
combinator library.

\vspace{1em}
\begin{definition}[Valid Message]\label{def:vvalid}
  A message $m$ is valid against a descriptor $d$, denoted $v \checkmark^d$ if
  both $v \checkmark_w$ and $v \checkmark_l^d$. 
\end{definition}

\vspace{1em}
\begin{definition}[Message Well-Formedness]\label{def:vwf}
  A message $m$ is well-formed, denoted $m \checkmark_w$, if all fields in
  the message have a unique field number, the fields are stored in sorted
  order, and all nested messages are also well-formed.
\end{definition}

\vspace{1em}
\begin{definition}[Legal Message]\label{def:vlegal}
  A message $m$ is legal with regards to descriptor $d$, denoted $m
  \checkmark_w^d$, if
  \begin{itemize}
  \item Totality, $\dom{m} = \dom{d}$.
  \item For all values in the message
    \begin{itemize}
    \item The presence of the field agrees with the cardinality in the
      descriptor.
    \item The type of the value agrees with the type of the field in the
      descriptor. For integer scalars, this condition also enforces the range
      bounds associated with the various Protobuf integer types (i.e.\ \ints{},
      \uintl{}, etc.) while for nested messages, that means $v' \checkmark^{d'}$.
    \end{itemize}
  \item At most one key in a \texttt{oneof} group defined in the descriptor is
    set. This is the only non-pairwise requirement that a well-formed message has
    to respect.
  \end{itemize}
\end{definition}

Descriptor legality (Definition~\ref{def:dlegal}), specifically the fact that
messages have to be \texttt{optional}, is important here since it makes it
possible to initialize a recursive message, and we can assume that anytime a
message is present, it's a non-default. This means that the Pollux check for if
a value is a default can be decidable, even though descriptors and fields don't
have a decidable equality instance defined.

\vspace{1em}
\begin{lemma}
  For all descriptors $d$, $\mathrm{init}(d) \checkmark^d$.
\end{lemma}

\vspace{1em}
\begin{definition}[Message Size]\label{def:vsize}
  The size of a message, $\mathrm{size}_v$, is defined below. 
  \begin{mathpar}
    \mathrm{size}_v(\llangle l \rrangle) = 1 + \mathrm{size}_v(l)

    \mathrm{size}_v(l) =
    \begin{cases}
      0 & \text{if } l = \texttt{[]} \\
      \mathrm{size}_v(v) + \mathrm{size}_v(l') & \text{if } l = (\_, v) :: l'
    \end{cases}

    \mathrm{size}_v(v) =
    \begin{cases}
      1 + \mathrm{size}_v(p) & \text{if } v = \impt{p} \\
      1 & \text{if } v = \optt{\text{none}} \\
      1 + \mathrm{size}_v(p) & \text{if } v = \optt{p} \\
      1 + \mathrm{size}_v(p_1) + \dots + \mathrm{size}_v(p_n) & \text{if } v = \rept{[p_1;\dots;p_n]}
    \end{cases}

    \mathrm{size}_v(p) =
    \begin{cases}
      \mathrm{size}_v(\llangle l' \rrangle) & \text{if } p = \llangle l'
                                              \rrangle \\
      1 & \text{otherwise}
    \end{cases}
  \end{mathpar}
\end{definition}

\subsection{The Transformation}\label{sec:proto-transform}

At the core of the Protobuf compatibility methodology is a transformation which
projects the view of a message under one descriptor to a related message under
another descriptor.

\subsubsection{Origins}

Using a function to transform the value isn't an arbitrary technical choice, nor
is it only useful to model potential descriptor changes to specific Protobuf
value. It arises naturally from the proofs of the intermediate format, and that
evolution is described here.

The first and strongest intermediate parser compatibility theorem is the
schema-correct one, Theorem~\ref{thm:inter-schema-correct}, where the relations
at play are strong enough that we can directly prove that the resulting value
from a serializer / parser round-trip under the same descriptor is equal. No
transformation is needed here.

However, the next step was to weak the compatibility antecedent to arrive at
Theorem~\ref{thm:inter-id-compat} where the descriptor is still the same but can
have no relation to the message being operated on. This is no longer possible
with our new totality conditions (Section~\ref{sec:value-totality}), but led to
the first example of a transformation function. The Intermediate ID
Compatibility relation has the \textsc{M-Drop-Unknown} rule, with no constrain
which value is dropped, and \textsc{M-Missing}, which can insert a new value
into the resulting message. Because of these rules, the induction hypothesis of
the theorem by itself is too weak to prove the whole theorem since the
compatibility relation has multiple $m_2$ related to the same $m_1$, $d_1$ and
$d_2$ so the proof must find a witness. Instead we use the transformation to
strengthen the inductive hypothesis into naming that witness and prove that
anything from the transformation will land inside the compatibility relation.

The final intermediate format theorem, Theorem~\ref{thm:inter-compat},
reinforces this. It uses the same technique with a much more complex
transformation function that actually mirrors the Protobuf transformation. Both
are a structural recursion on the reader's descriptor field list, building the
new interpretation of the value under the new descriptor using the writer's
descriptor. The Protobuf one also has the stronger justification of enforcing
the totality requirement.

\subsubsection{Definition}

% Local Variables:
% citar-bibliography: ("../pollux.bib")
% TeX-master: "../pollux.tex"
% TeX-command-default: "Make"
% jinx-local-words: "protobuf"
% End:
